% Transcription — fonds Alexandre Grothendieck, shelfmark 33, batch 9 (pages 161-180).
% Established from the facsimile archives/batches/33/batch-09.pdf, per the skill
% transcribe-grothendieck. Pages 162, 167, 170, 172, 174, 176, 178 and 180 are
% unrelated typescripts (Bourbaki « tribu », papers in English), page 168 is blank,
% page 163 is a cover whose words become the section heading — all absent.
% Pass: Opus 5.5 (claude-opus-5-5), 2026-10-03 — first pass, unchecked against the pages by a human.
\documentclass[11pt,a4paper]{article}
\input{../preamble/grothendieck}

\folder{33}
\batch{9}
\pages{161}{180}
\dating{1965-[vers 1971]}
\watermark{Édition de démonstration}
\foldertitle{SGA 1965. MS [Manuscrit] brouillon : notes manuscrites (s.d.).}

\begin{document}

\page{161}
\note{la page continue un argument commencé avant ce lot : les classes $\overline{\alpha^{0}}$, $\overline{\beta^{0}}$, les correspondances $T$, $T'$ et l'accouplement $\langle \xi^{0}, \eta^{0}\rangle$ sont introduits plus haut, hors du lot.}

\drawing{en marge gauche, le diagramme $T \times_{Z} T' \to T, T' \to X \times_{S} Y = Z$ :}

\begin{tikzcd}[column sep=small, row sep=small]
 & T \times_{Z} T' \arrow[dl] \arrow[dr] & \\
T \arrow[dr] & & T' \arrow[dl] \\
 & X \times_{S} Y = Z &
\end{tikzcd}

À vrai dire, \struck{\ill{}} si $\overline{T}$ est l'image de $T$ dans
$Z = X \times_{S} Y$, on veut que $\overline{\alpha^{0}}$ « vive » dans
\struck{\uncertain{l'image}}
\[
  R^{0}\Gamma_{\overline{T}}\bigl(R\mathcal{H}om(\mathrm{pr}_{1}^{*}(F), \mathrm{pr}_{2}^{!}(G))\bigr)
  = R^{0}\Gamma_{X \times_{S} Y}\bigl(R\underline{\Gamma}_{\overline{T}}\, R\mathcal{H}om( \ , \ )\bigr)
\]
et de même que $\overline{\beta^{0}}$ « vive » \struck{\uncertain{dans l'image}} dans
\[
  R^{0}\Gamma_{\overline{T'}}\bigl(R\mathcal{H}om(\mathrm{pr}_{2}^{*}(G), \mathrm{pr}_{1}^{!}(F))\bigr)
  = R^{0}\Gamma_{X \times_{S} Y}\bigl(R\underline{\Gamma}_{\overline{T'}}, R\mathcal{H}om \ldots\bigr)
\]

Ainsi, on peut définir leur cup-produit
\[
  \overline{\alpha^{0}} \cup \overline{\beta^{0}}
  \in R^{0}\Gamma_{\overline{T} \cap \overline{T'}}\, \underline{K}^{\bullet}_{X \times_{S} Y}
  \simeq R^{0}\Gamma_{\overline{T} \cap \overline{T'}}\, K^{\bullet}_{\overline{T} \times_{S} \overline{T'}}
\]
\note{l'indice du dernier $K^{\bullet}$ se lit $\overline{T} \times_{S} \overline{T'}$ ; la lecture de l'indice $S$ est incertaine.}

Cela permet donc de définir, si
$\overline{T} \cap \overline{T'} = \bigcup_{\alpha} S_{\alpha}$ \struck{réunion}
les $S_{\alpha}$ ouverts deux à deux disjoints (en nb fini), de définir des
composantes canoniques $(\overline{\alpha^{0}} \cup \overline{\beta^{0}})_{\alpha}$
\struck{\ill{}} et de même \struck{\ill{}} $\langle \xi^{0}, \eta^{0}\rangle_{S_{\alpha}}$
telles que
\[
  \langle \xi^{0}, \eta^{0}\rangle = \sum_{\alpha} \langle \xi^{0}, \eta^{0}\rangle_{S_{\alpha}} .
\]

\emph{Exemple.} Supposons $\overline{T} \cap \overline{T'} = \emptyset$. Alors
le nb de Lefschetz est nul.

\section{Compléments dualité}
\note{titre pris de la page 163, une couverture qui ne porte que ces mots, de sa main, en haut à droite.}

\page{164}
\marginal{en haut à droite, au crayon : codim inj $F \leq 2 \dim$ \uncertain{$\operatorname{Supp}$} $F + \operatorname{Sup}_{x \in X} \dim \operatorname{inj} F_{\bar{x}}$}

\emph{Appendices possibles}

(1) Majoration de la codim. inj. via la dim. inj. des fibres. Dim. coh.
stricte d'un préschéma (en particulier d'un pr. alg.)

(Consulter notes de Verdier)
\marginal{au crayon, sous la ligne : $\text{c-dim-inj}(F) \leq \operatorname{Sup}_{x \in X,\, F_{\bar{x}} \neq 0} \bigl[2 \dim \mathcal{O}_{X,x} + \dim \operatorname{inj} F_{\bar{x}}\bigr]$}

(2) Cas d'un Anneau $A$ de torsion loc. ct et premier aux car. résiduelles,
et $R^{\bullet}_{X}$ complexe de faisceaux de $A_{X}$-modules loc. cts,
dualisant fibre par fibre (sur $X$ régulier …)

\emph{NB.} On peut globaliser, et considérer des schémas \uncertain{relatifs}
sur $(X_{\text{ét}}, (\mathbb{Z}/n\mathbb{Z})_{X})$ …

\struck{(3) Applications aux ths. de Lefschetz cohomologiques locaux et globaux
[Pourrait donner \uncertain{un} article séparé …]}
\note{l'alinéa (3) est barré de traits obliques.}

\page{165}
\emph{Sujets de recherches}
\begin{enumerate}
  \item[a)] $\Lambda^{i}$ dans catégories dérivées, $\lambda$-anneaux et structures annexes
  \item[b)] Dualité analytique
  \item[c)] Descente cohomologique et applications (finitude analytique et \ill{} …)
  \item[d)] Künneth pour complexes quasi-cohérents
  \item[e)] Dualité pour les complexes d'op. différentiels, et site stratifiant …
  \item[f)] \struck{Profondeurs cohomologiques et homotopiques, et ths de Lefschetz locaux et globaux.}
\end{enumerate}
\note{une accolade réunit b) et c), une autre d) et e) ; f), sous un trait horizontal, est barré de traits obliques.}

\page{166}
\emph{Validité des théorèmes de dualité}

1) \emph{Dualité globale}
\[
  \operatorname{Hom}(F_{\bullet}, R^{!}f(G^{\bullet}))
  \xleftarrow{\ \sim\ } \operatorname{Hom}(R_{!}f(F_{\bullet}), G^{\bullet})
\]
$f \colon (X, A_{X}) \to (Y, A_{Y})$ morphisme compactifiable et
\uncertain{lissifiable} de préschémas, annelés par $A_{X}$ et $A_{Y}$,
\struck{\ill{}} faisceaux d'anneaux \struck{de torsion} tels que pour tt pt
géom. $\bar{y}$ de $Y$, $\exists\, n_{\bar{y}}$ tel que
$n_{\bar{y}} A_{Y,\bar{y}} = 0$ \struck{et un anneau} \uncertain{a} $n_{\bar{y}}$
premier à $\operatorname{car} k(\bar{y})$.

2) \emph{Dualité locale}
\[
  F_{\bullet} \xrightarrow{\ \sim\ } D_{X} D_{X}(F_{\bullet})
\]
isom., $D_{X}(F_{\bullet})$ à coh. constructible,

$X$ préschéma \struck{loc} noeth. régulier, annelé par $A_{X}$ faisceau
d'anneaux loc. ct à fibres \struck{noethériens} \ill{} des anneaux de torsion
noethériens, premiers aux car. résiduelles, $R^{\bullet}_{X}$ un
\add{complexe de} faisceaux de $A_{X}$-Modules loc. constants, et dualisant
\add{sur les} fibres \struck{\uncertain{géométriques}}
\[
  D_{X}(F_{\bullet}) = R\mathcal{H}om(F_{\bullet}, A_{X}),
  \qquad F_{\bullet} \in \operatorname{Ob} D^{b}(X)
\]
\marginal{dans un cadre en bas à droite : + résolution et pureté sur $X$ …}

\page{169}
\emph{Énoncés plausibles} (sur préschéma excellent de dim finie).
\[
  D_{X} \
  \begin{cases}
    D_{c}(X)^{\circ} \xrightarrow{\ \approx\ } D_{c}(X) \\
    D_{c}^{-}(X)^{\circ} \xrightarrow{\ \approx\ } D_{c}^{+}(X) \\
    D_{c}^{+}(X)^{\circ} \xrightarrow{\ \approx\ } D_{c}^{-}(X) \\
    D_{c}^{b}(X) \xrightarrow{\ \approx\ } D_{c}^{b}(X)
  \end{cases}
\]

$F_{\bullet} \in \operatorname{Ob} D_{c}^{-}(X)$,
\add{$G^{\bullet} = D_{X}(F_{\bullet}) \in \operatorname{Ob} D_{c}^{+}(X)$} je dis que

(*) $F_{\bullet}$ de Tor-dim finie $\Longleftrightarrow$ $G^{\bullet}$ de
pseudo-dim. inj. finie [i.e. $\exists\, N$ tel que pour tt \struck{$G$} $P$
\struck{$\in \operatorname{Ob} D_{c}$} \add{\uncertain{parfait}} constructible,
$\mathcal{E}xt^{i}(P, G^{\bullet}) = 0$ pour $i > N$]

(**) $\Longleftrightarrow$ fibres \struck{\ill{}} $G^{\bullet}_{\bar{x}}$ sont de
dim. inj. finie (sur les $A$-Modules)

\emph{N.B.} Si $X$ est de dim. coh. finie, on sait que $G^{\bullet}$ donné ;
$G^{\bullet}$ pseudo-dim. inj. finie $\Longleftrightarrow$ $G^{\bullet}$ de dim.
injective finie.
\note{un trait vertical au crayon borde ce N.B. à gauche.}

Dém. de (*). L'implication $\Rightarrow$ formelle, en utilisant que
$R^{\bullet}_{X}$ est de pseudo-dim. injective finie, et
\[
  R\mathcal{H}om(P_{\bullet}, D_{X}(F_{\bullet}))
  \simeq R\mathcal{H}om(P_{\bullet}, R\mathcal{H}om(F_{\bullet}, R^{\bullet}_{X}))
  \simeq R\mathcal{H}om(P_{\bullet} \overset{L}{\otimes} F_{\bullet}, R^{\bullet}_{X})
\]
Pour l'implication inverse, on utilise la bidualité pour
$F_{\bullet} \overset{L}{\otimes} P$ …

\emph{N.B.} Si les fibres \ill{} $A_{\bar{x}}$ sont Gorenstein, alors en
admettant l'équivalence (**), on veut que si
$F_{\bullet} \in \operatorname{Ob} D_{c}(X)$, Tor-dim finie $\Longleftrightarrow$
pseudo-dim inj. finie, et $D_{X}$ induit une anti-autoéquivalence dans la
catégorie desdits complexes …

\page{171}
\marginal{en diagonale, en haut à gauche : Généraliser : aux faisceaux d'anneaux \ill{} finis \ill{}}

\struck{\ill{}} \emph{à vérifier !}

\emph{Théorème.} Soit $X$ \add{lisse} de type fini sur le corps $k$,
$\dim X \leq n$, $A = \mathbb{Z}/N\mathbb{Z}$, $N$ premier à car. $k$,
\struck{$A$} $F$ et $G$ deux $A$-\struck{\ill{}} Modules. Supposons que $F$ soit
plat \uncertain{$=$} $G$ : fibres injectives (= plates, dans le cas particulier
envisagé !). Alors

(i) Si $F$ constructible, alors $\operatorname{Ext}^{i}(F, G) = 0$ pour $i > 2n$.

(ii) Si $k$ est séparablement clos, \uncertain{alors} \struck{pareillement}
\[
  \operatorname{Ext}^{i}(\cdot\,, G) = 0 \quad \text{pour } i > n, \qquad
  \operatorname{Ext}^{i}(F, G) = 0 \quad \text{pour } i > 2n .
\]
\note{le premier argument du premier $\operatorname{Ext}^{i}$ est surchargé (un $X$ ou un $A$ corrigé) ; la borne « $i > n$ » est d'une lecture incertaine.}

\emph{Dém.} (i) se ramène par \uncertain{passage} à la limite à (ii). Dans
(ii) on peut supposer $k$ alg. clos, et \struck{un argument général} il suffit
de prouver \add{\uncertain{supposons} d'abord $G$ : fibres injectives.} la
\uncertain{première} relation. \add{Un argument général de \uncertain{récurrence}
sur \add{la formule} longueur $F$} et \ill{} : on simplifie \ill{}, partant
d'un générateur de la forme $i_{!}(A_{X'})$, où $X'$ affine et $X' \to X$
étale, on \ill{} d'un \ill{} \uncertain{ouvert affine} $U$ de $X$. Car on se
ramène au cas où $X = U$ est aussi affine [car on peut remplacer $X$ par $U$,
$G$ par $G|U$]
\note{l'ordre des insertions interlinéaires de ce paragraphe, séparées du texte par deux traits obliques, est incertain ; la page 173 en donne la suite.}

\page{173}
\note{suite de la page 171 (la page 172 est le verso d'un tapuscrit étranger).}
et $F$ constructible. De plus, par passage à la limite, on peut supposer $G$
\ill{} constructible [Attention, est-il vrai que $G$ est limite ind.
\add{filtrante} des \struck{\ill{}} faisceaux const. à fibres
\emph{injectives} ??].
\marginal{un « ? » en marge devant « filtrante »}

\page{175}
\note{page de calculs épars sur des suites spectrales ; on transcrit les formules dans l'ordre de la page et on décrit les figures.}
\[
  \mathcal{E}xt^{p}(F, G)
\]
\[
  \operatorname{Ext}^{n}(i_{!}(F), G) \simeq \operatorname{Ext}^{n}(F, Ri^{!}(G))
\]
\[
  \operatorname{Ext}^{*}(F, Ri^{!}(G)) \Longleftarrow E_{2}^{pq}
  = \operatorname{Ext}^{p}(F, R^{q}i^{!}(G))
\]
$Y \hookrightarrow X$
\[
  \operatorname{Ext}^{p}(F, \underline{R^{q}i^{!}G})
\]
\drawing{un triangle rectangle, l'hypoténuse redoublée.}

$\mathcal{E}xt^{p}$, avec au-dessus les indices \uncertain{$2d$, $2d-1$, $2d-2$, …, $0$} :
\[
  \mathcal{E}xt^{0}, \mathcal{E}xt^{1}\ \mathcal{E}xt^{2} \,/\,
  \mathcal{E}xt^{3}\ \mathcal{E}xt^{4} \,/ \quad \cdots \quad /\,
  \mathcal{E}xt^{2d-1}, \mathcal{E}xt^{2d}
\]
\[
  H^{p}(X, \mathcal{E}xt^{q}) \qquad
  \text{\struck{$H^{q}$}}\ \underline{H}^{n}(X, \mathcal{E}xt^{0}),\ H^{n-1}\mathcal{E}xt^{1},
\]
$p \leq$
\[
  n \leq 2d \qquad
  n-1 \leq 2(d-1) \qquad
  n-3 \leq 2(d-2)
\]
\[
  \text{\struck{$n-1 \leq 2d-1$}} \qquad
  n-2 \leq 2(d-1) \qquad
  n-4 \leq 2(d-2) \qquad
  \underline{2d-1}
\]
\drawing{deux axes, le long de l'axe horizontal un segment de $0$ à $2d$ marqué de points ; puis un triangle sous une droite de pente $-1$, strié de traits horizontaux, l'abscisse $2d$ marquée.}
\[
  \text{\struck{$\operatorname{Ext}^{*}(F, G) \simeq$}}
\]
\[
  R\mathcal{H}om(F, \text{\struck{$D$}}G) \simeq R\mathcal{H}om(F \otimes DG, K_{X})
\]
\[
  \operatorname{Ext}^{n}(X; F, G) \Longleftarrow E_{2}^{pq}
  = \operatorname{Ext}^{p}(X; \mathcal{E}xt^{q}(X; G, A_{X}), A_{X})
\]
$\operatorname{Ext}^{2d}(X; F, G)$
\drawing{un rectangle dont le coin supérieur droit est marqué $2d$ ; deux axes $(p, q)$, des points sur l'axe horizontal jusqu'à $2d$.}
\[
  \operatorname{Ext}^{*}(X; F, A_{X}) \Longleftarrow H^{p}(X, \mathcal{E}xt^{q}(F, A_{X}))
  \qquad 2d
\]
$H^{2d}(X,$
\drawing{un segment marqué $2d$ à son extrémité droite.}

\page{177}
\note{nouvelle page de calculs, de la même facture que la page 175.}
\[
  \mathcal{E}xt^{i}(F, G) \qquad \text{\struck{$F$}}\ DG
\]
\[
  R\mathcal{H}om(F, G) \simeq R\mathcal{H}om(F, DDG)
  \simeq R\mathcal{H}om(F \overset{L}{\otimes} DG, R^{\bullet}_{X})
\]
\[
  F \otimes DG \qquad -2d \longrightarrow 0 \qquad -2d \longrightarrow 0
\]
$\mathcal{E}xt^{i}$ \quad $\mathcal{H}_{i}(F \otimes DG)$ de dim $\leq i$
\[
  D(F \overset{L}{\otimes} DG \qquad \mathcal{E}xt^{i}(F, G)
\]
Encadré :
\[
  \mathcal{E}xt^{*}(F, G) \Longleftarrow E_{2}^{pq}
  = \mathcal{E}xt^{p}(F \otimes \mathcal{E}xt^{q}(G, A_{X}), A_{X})
\]
\note{l'indice inférieur de $E^{pq}$ se lit mal ($1$ ou $2$) ; l'exposant de $\mathcal{E}xt^{p}$ porte une rature, et celui de $\mathcal{E}xt^{q}$ semble surchargé en $-q$.}

\drawing{axes $(p, q)$ ; un rectangle de sommets $0$ et $2d$, sa diagonale fortement tracée, la droite $p + q = 0$ prolongée au-delà ; au-dessus, une courte inscription biffée, \uncertain{$q^{4} = q \geq q$}.}

$F \otimes \mathcal{E}xt^{2d}(G, A)$ \quad \struck{$\mathcal{E}xt^{p}($}

\uncertain{supp} $\mathcal{E}xt^{-q}(G, A_{X})$ de \struck{\ill{}} \uncertain{codim.} $\geq -q$
\struck{de dim. $\leq$}, donc $\mathcal{E}xt^{p}( \ , A_{X}) = 0$ si $p < -q$, i.e. $p + q < 0$.
\[
  \text{\struck{$\mathcal{E}xt^{*}$}}\ \mathcal{E}xt^{2d}(F, G)
  \simeq \mathcal{E}xt^{2d}(F \otimes \mathcal{H}om(G, A_{X}), A_{X})
\]

\page{179}
\note{page de calculs ; la suite exacte de cohomologie locale et sa duale sont mises en regard, terme à terme, par des croix.}
\[
  \cdots \to R^{i}\Gamma_{x}(F) \to R^{i}\Gamma_{X}(F) \to R^{i}\Gamma_{U}(F)
  \to R^{i+1}\Gamma_{x}(F) \to \cdots
\]
\[
  R^{-i}\Gamma_{X}(DF) \leftarrow R^{-i}\Gamma_{x}DF \leftarrow R^{-i-1}\Gamma_{U}(DF)
  \leftarrow R^{-i-1}\Gamma_{x}DF \leftarrow \cdots
\]
\drawing{sous trois des termes de la seconde ligne, des flèches verticales vers $A$ ; sous la première, l'annotation « dualité parfaite triviale (via \uncertain{pureté} …) » ; sous les deux autres, réunies par des arcs, « dualité parfaite ?? ».}

$f^{!}R\mathcal{H}om(F^{\bullet}, G^{\bullet})$ \qquad $f^{!}D = Df^{*}$

\[
  R\mathcal{H}om(F, G) \qquad D(F \otimes DG) \qquad
  R\mathcal{H}om(G, DF) \simeq D(G \overset{L}{\otimes} F)
\]
\[
  R\mathcal{H}om(F, D(DG)) \qquad
  D(F \otimes DG) \simeq DR\mathcal{H}om(F, G)
\]
\note{dans la dernière formule, un $D$ est ajouté sous $G$ ; la lecture de cette correction est incertaine.}
\[
  D(\mathcal{H}^{i}_{x}(F)) \longrightarrow \text{\struck{$D$}}\ \mathcal{H}^{-i}_{x}(DF)_{x}
  \qquad F \to DDF
\]
\[
  \text{\struck{$\mathcal{H}^{i}_{x}(F) \to$}} \qquad
  \mathcal{H}^{i}(F)_{x} \longrightarrow D_{x}\,\underline{\mathcal{H}^{-i}_{x}(DF)}
  \qquad D_{x} \text{\struck{$\mathcal{H}$}}\, D_{x}(\mathcal{H}^{i}(F)_{x})
\]
\drawing{des intervalles $]\ \ [$ et $]\!-\!-\!-\![$ ; un axe portant $M$, $0$, $N$ ; une barre de longueur $-2n$ ; en marge gauche $F \otimes DG$, $-2n, 0$.}

Tor-dim. $F_{\bullet}$ \uncertain{concentrée} dans $[N, M]$
\[
  \Updownarrow
\]
\uncertain{loc.} inj.-dim. $D(F_{\bullet})$ \uncertain{concentrée} dans $[-2n + N, M]$

Inj.-dim. $F$ \uncertain{concentrée} dans $[P, Q]$
\[
  \Updownarrow
\]
Tor-dim. $DF$ \uncertain{concentrée} dans $[P - 2n, Q]$
\[
  DF \overset{L}{\otimes} G \simeq DR\mathcal{H}om(G, F),
\]
\[
  \text{\struck{$DF \otimes DDG \simeq D(F \otimes D$}}
  \qquad [P, Q]
\]
\note{les bornes $[-2n + N, M]$ et $[P - 2n, Q]$ sont transcrites telles qu'écrites ; l'asymétrie entre les deux équivalences n'est pas commentée sur la page.}

\end{document}
